\documentclass[12pt, letterpaper]{article}

\usepackage[utf8]{inputenc}
\usepackage[T1]{fontenc}
\usepackage[spanish]{babel}
\usepackage{geometry}
\usepackage{amsmath, amssymb, amsthm}
\usepackage{booktabs}
\usepackage{array}
\usepackage{microtype}
\usepackage{setspace}
\usepackage{parskip}
\usepackage{titlesec}
\usepackage{fancyhdr}
\usepackage{enumitem}
\usepackage{bm}
\usepackage{mathtools}
\usepackage{tcolorbox}
\usepackage{hyperref}

\tcbuselibrary{skins, breakable}

\geometry{top=2.8cm, bottom=3.0cm, left=3.2cm, right=3.2cm, headheight=14pt}
\setstretch{1.20}

\pagestyle{fancy}
\fancyhf{}
\fancyhead[C]{\small\textsc{Econometría --- Gujarati --- Ejercicio 13.14}}
\fancyfoot[C]{\small\thepage}
\renewcommand{\headrulewidth}{0.4pt}
\renewcommand{\footrulewidth}{0pt}

\titleformat{\section}
  {\large\bfseries\scshape}{\thesection.}{0.6em}{}
  [\vspace{-4pt}\rule{\linewidth}{0.4pt}]
\titleformat{\subsection}{\normalsize\bfseries}{\thesubsection.}{0.5em}{}
\titleformat{\subsubsection}{\normalsize\itshape}{}{0em}{}

\newtcolorbox{clave}{
  colback=white, colframe=black,
  boxrule=0.5pt, arc=3pt,
  left=8pt, right=8pt, top=6pt, bottom=6pt,
  breakable
}

\newcommand{\E}{\operatorname{E}}
\newcommand{\Cov}{\operatorname{Cov}}
\newcommand{\Var}{\operatorname{Var}}

\begin{document}

\begin{center}
  {\LARGE\bfseries Evaluación de la Advertencia de Henry Theil}\\[6pt]
  {\LARGE\bfseries sobre la Degradación de la Inferencia Estadística}\\[12pt]
  {\large\itshape Minería de datos, inflación del error Tipo~I}\\
  {\large\itshape y colapso de los intervalos de confianza nominales}\\[14pt]
  \rule{0.55\linewidth}{0.6pt}\\[8pt]
  {\normalsize Ejercicio~13.14 --- \textit{Econometría}, Gujarati \& Porter, 5.a ed.}\\[4pt]
  {\small\textsc{Análisis conceptual, teórico y algebraico}}
\end{center}

\vspace{10pt}

%------------------------------------------------------------------------------
\section{Planteamiento: la afirmación de Theil}
%------------------------------------------------------------------------------

Henry Theil advierte que cuando los intervalos de confianza y los
estadísticos de prueba se calculan a partir de la \emph{regresión final}
de una estrategia de búsqueda de especificación, sus propiedades nominales
ya no son válidas. Concretamente señala que:

\begin{itemize}[leftmargin=1.6cm]
  \item Un coeficiente reportado con un \textbf{intervalo de confianza del
    $95\,\%$} puede, en realidad, tener una cobertura de apenas el
    $\mathbf{80\,\%}$.
  \item Un nivel de significancia nominal del $\mathbf{1\,\%}$ puede
    corresponder, en la práctica, a un nivel real del $\mathbf{10\,\%}$.
\end{itemize}

La afirmación no es una opinión metodológica vaga: tiene un sustento
probabilístico preciso que se puede derivar algebraicamente a partir de
la estructura del proceso de selección de modelos. El objetivo de esta
evaluación es demostrar dicho sustento, identificar sus supuestos y
discutir sus implicaciones para la práctica econométrica.

%------------------------------------------------------------------------------
\section{Fundamento conceptual: inferencia clásica y sus supuestos implícitos}
%------------------------------------------------------------------------------

\subsection{El supuesto de pre-especificación}

La validez de las pruebas $t$, $F$ y de los intervalos de confianza
estándar descansa en un supuesto que raramente se menciona de forma
explícita: el modelo fue \textbf{especificado antes de observar los datos},
o al menos antes de examinar las correlaciones entre las variables de
interés. Bajo esta condición, la probabilidad de rechazar una hipótesis
nula verdadera es exactamente $\alpha$ (el nivel de significancia elegido),
y el intervalo de confianza al $(1-\alpha)\times 100\,\%$ contiene el
verdadero parámetro con esa misma frecuencia en muestreos repetidos.

\subsection{La estrategia de regresión y el sesgo de selección}

En la práctica, los investigadores suelen seguir una \textbf{estrategia de
regresión}: se prueban múltiples combinaciones de variables, formas
funcionales y submuestras hasta que el modelo resultante luce
``estadísticamente satisfactorio'' (coeficientes significativos, signos
coherentes con la teoría, $R^2$ elevado). La regresión \emph{final} que
se publica es la sobreviviente de un proceso de selección.

Este proceso introduce un \textbf{sesgo de selección de modelos}: los
datos han sido examinados repetidamente y el modelo reportado es,
por construcción, el que mejor se ajusta a esa muestra particular.
Los estadísticos de ese modelo están sesgados hacia la significancia,
no porque el efecto sea real, sino porque el proceso de búsqueda garantiza
que solo los modelos con resultados ``favorables'' llegan a publicarse.
En términos coloquiales: los datos han sido ``torturados hasta confesar''.

\subsection{Grados de libertad implícitos no contabilizados}

Cada vez que el investigador examina una especificación alternativa y la
descarta, consume información de la muestra sin que ello quede reflejado
en los grados de libertad del modelo final. El resultado es que los
errores estándar reportados \emph{subestiman} la verdadera variabilidad
del estimador, y los intervalos de confianza son \emph{más estrechos}
de lo que deberían ser. Este es el mecanismo por el cual un intervalo
nominal del $95\,\%$ ofrece cobertura real inferior.

%------------------------------------------------------------------------------
\section{Fundamento algebraico: la fórmula de Lovell}
%------------------------------------------------------------------------------

\subsection{Derivación del nivel de significancia real}

Sea $c$ el número total de variables candidatas examinadas durante el
proceso de búsqueda y $k$ el número de variables retenidas en el modelo
final, con $k \leq c$. Supóngase que cada prueba individual se realiza
al nivel nominal $\alpha$, y que las pruebas son aproximadamente
independientes. La probabilidad de que, bajo la hipótesis nula global,
\emph{al menos una} de las $c$ pruebas individuales resulte significativa
es:
\begin{equation}
  P(\text{al menos un rechazo erróneo}) = 1 - (1-\alpha)^c.
  \label{eq:bonferroni_c}
\end{equation}
Sin embargo, el investigador no reporta todas las pruebas realizadas,
sino solo las $k$ variables seleccionadas. El nivel de significancia
efectivo del modelo final, ajustado por la búsqueda de las mejores $k$
variables entre $c$ candidatas, es la \textbf{fórmula de Lovell}:
\begin{equation}
  \alpha^* = 1 - (1 - \alpha)^{c/k}.
  \label{eq:lovell_exacta}
\end{equation}
Para $\alpha$ pequeño, la expansión de Taylor de primer orden de
$(1-\alpha)^{c/k} \approx 1 - (c/k)\alpha$ produce la aproximación
lineal:
\begin{equation}
  \alpha^* \approx \frac{c}{k}\,\alpha.
  \label{eq:lovell_aprox}
\end{equation}

\subsection{Verificación numérica del ejemplo de Theil: $1\,\% \to 10\,\%$}

Theil afirma que un nivel de significancia nominal del $1\,\%$ puede
corresponder en realidad a un $10\,\%$. Este resultado se reproduce
exactamente con $c = 20$ variables candidatas y $k = 2$ variables
retenidas:
\begin{equation}
  \alpha^*
  \approx \frac{c}{k}\,\alpha
  = \frac{20}{2} \times 0{,}01
  = 10 \times 0{,}01
  = 0{,}10.
  \label{eq:ejemplo_1pct}
\end{equation}
Usando la fórmula exacta~\eqref{eq:lovell_exacta}:
\begin{equation}
  \alpha^* = 1 - (1 - 0{,}01)^{20/2}
           = 1 - (0{,}99)^{10}
           = 1 - 0{,}9044
           = 0{,}0956
           \approx 9{,}6\,\%.
  \label{eq:ejemplo_1pct_exacto}
\end{equation}
Ambas expresiones confirman el orden de magnitud: el nivel real es
aproximadamente diez veces mayor que el nominal reportado.

\subsection{Verificación numérica del ejemplo de Theil: $95\,\% \to 80\,\%$}

El nivel de confianza nominal es $1 - \alpha = 0{,}95$, es decir,
$\alpha = 0{,}05$. Supóngase que el proceso de búsqueda involucra
$c = 4$ variables candidatas con $k = 1$ variable retenida:
\begin{equation}
  \alpha^*
  \approx \frac{c}{k}\,\alpha
  = \frac{4}{1} \times 0{,}05
  = 0{,}20.
  \label{eq:ejemplo_5pct}
\end{equation}
El nivel de confianza real es entonces:
\begin{equation}
  1 - \alpha^* = 1 - 0{,}20 = 0{,}80 = 80\,\%.
  \label{eq:confianza_real}
\end{equation}
Verificación con la fórmula exacta:
\begin{equation}
  \alpha^* = 1 - (1 - 0{,}05)^{4/1}
           = 1 - (0{,}95)^{4}
           = 1 - 0{,}8145
           = 0{,}1855
           \approx 18{,}6\,\%.
  \label{eq:ejemplo_5pct_exacto}
\end{equation}
El nivel de confianza real es $1 - 0{,}186 \approx 81{,}4\,\%$,
consistente con la cota del $80\,\%$ señalada por Theil. Las dos
afirmaciones del enunciado quedan así demostradas algebraicamente.

\subsection{Resumen de los resultados numéricos}

\begin{center}
\renewcommand{\arraystretch}{1.4}
\small
\begin{tabular}{cccccc}
\toprule
\textbf{Nivel nominal} & $c$ & $k$ & $c/k$ & $\alpha^*$ (aprox.) & \textbf{Nivel real} \\
\midrule
$\alpha = 0{,}01\;(1\,\%)$  & 20 & 2 & 10 & $0{,}10$  & $\approx 10\,\%$ \\
$\alpha = 0{,}05\;(5\,\%)$  & 4  & 1 & 4  & $0{,}20$  & $\approx 80\,\%$ de confianza \\
$\alpha = 0{,}05\;(5\,\%)$  & 3  & 1 & 3  & $0{,}15$  & $\approx 85\,\%$ de confianza \\
\bottomrule
\end{tabular}
\end{center}

%------------------------------------------------------------------------------
\section{Evaluación crítica de la afirmación de Theil}
%------------------------------------------------------------------------------

\subsection{Aspectos en los que la afirmación es correcta}

\begin{enumerate}[leftmargin=1.8cm, label=\textbf{\arabic*.}]

  \item \textbf{Validez matemática.} La fórmula de Lovell
    [\eqref{eq:lovell_exacta}--\eqref{eq:lovell_aprox}] es un resultado
    probabilístico riguroso, no una intuición. La degradación de los
    niveles nominales bajo búsqueda de especificación es un hecho
    demostrable, no una opinión.

  \item \textbf{Prevalencia del problema en la práctica.} La minería
    de datos no es una anomalía rara: es el procedimiento habitual
    en econometría aplicada. La mayoría de las publicaciones empíricas
    no reportan el número de especificaciones exploradas antes del
    modelo final, lo que hace imposible para el lector calibrar
    $\alpha^*$.

  \item \textbf{La recomendación de interpretar ``liberalmente'' es
    operativamente correcta.} En ausencia de información sobre el
    proceso de búsqueda, la única respuesta racional del lector es
    aplicar un descuento al nivel de significancia reportado. Theil
    formaliza este descuento de manera que puede cuantificarse.

\end{enumerate}

\subsection{Limitaciones y matices}

\begin{enumerate}[leftmargin=1.8cm, label=\textbf{\arabic*.}]

  \item \textbf{La fórmula supone pruebas independientes.}
    La derivación de~\eqref{eq:lovell_exacta} asume que las $c$ pruebas
    son aproximadamente independientes. En la práctica, las variables
    candidatas suelen estar correlacionadas, lo que hace que la inflación
    real de $\alpha^*$ sea menor que la predicha por Lovell. La fórmula
    es por tanto una \emph{cota superior} del nivel real, no una estimación
    exacta.

  \item \textbf{No toda búsqueda de especificación es invalidante.}
    Existen procedimientos de selección de variables con propiedades
    estadísticas controladas: el criterio de información de Akaike (AIC),
    el criterio bayesiano de Schwarz (BIC) y los métodos LASSO penalizan
    la complejidad del modelo de forma que la tasa de falsos positivos
    puede mantenerse bajo control. Theil escribe en un contexto anterior
    al desarrollo masivo de estos métodos.

  \item \textbf{La solución propuesta (interpretación ``liberal'') es
    paliativa, no estructural.} Ajustar subjetivamente los niveles de
    confianza reportados no sustituye la necesidad de transparencia en
    el proceso de modelado. La solución correcta es reportar todas las
    especificaciones exploradas, aplicar la corrección de Lovell
    explícitamente, o utilizar métodos de selección con garantías formales.

\end{enumerate}

%------------------------------------------------------------------------------
\section{Implicaciones para la práctica econométrica}
%------------------------------------------------------------------------------

La advertencia de Theil impone cuatro obligaciones metodológicas concretas:

\begin{enumerate}[leftmargin=1.8cm, label=\textbf{P\arabic*.}]

  \item \textbf{Transparencia en el proceso de búsqueda.} El investigador
    debe reportar $c$ y $k$, de modo que el lector pueda aplicar
    la corrección~\eqref{eq:lovell_aprox} y obtener $\alpha^*$.

  \item \textbf{Análisis de sensibilidad.} Un resultado solo merece
    confianza si los coeficientes clave son estables ante variaciones
    razonables en la especificación (inclusión o exclusión de controles,
    transformaciones funcionales, submuestras).

  \item \textbf{Uso de criterios de información penalizados.} El AIC y
    el BIC introducen una penalización por el número de parámetros que
    mitiga parcialmente el sobreajuste. No eliminan el problema de la
    minería de datos, pero reducen su severidad.

  \item \textbf{Distinción entre significancia estadística y relevancia
    económica.} Dado que la significancia nominal es una cota inferior
    de la significancia real, un coeficiente borderline ($p \approx 0{,}05$)
    tras búsqueda de especificación puede tener $\alpha^* > 0{,}20$.
    La pregunta que debe primar es si la magnitud del efecto es
    sustantiva desde el punto de vista económico o de política pública.

\end{enumerate}

\begin{clave}
  \textbf{Conclusión:} La afirmación de Theil es técnicamente correcta y
  empíricamente relevante. La fórmula de Lovell demuestra algebraicamente
  que la búsqueda de especificación infla el error Tipo~I en un factor de
  $c/k$, degradando los intervalos de confianza y los niveles de
  significancia de manera proporcional a la intensidad del proceso de
  búsqueda. Sus limitaciones son que la fórmula es una cota superior
  (las pruebas reales no son independientes) y que la recomendación
  de interpretar ``liberalmente'' es una respuesta paliativa que no
  reemplaza la transparencia metodológica. El legado de Theil es haber
  cuantificado el costo estadístico de la minería de datos, convirtiendo
  una crítica cualitativa en un resultado numérico verificable.
\end{clave}

\end{document}
